\begin{lemma}
\label{lemma-base-change-q-injective}
With $f : X \to S$ and $n$ as in Remark \ref{remark-base-change-holds}
assume for some $q \geq 1$ we have $BC(f, n, q - 1)$. Then
for every commutative diagram
$$
\xymatrix{
X \ar[d]_f & X' \ar[l] \ar[d]_{f'} & Y \ar[l]^h \ar[d]^e \\
S & S' \ar[l] & T \ar[l]_g
}
$$
with $X' = X \times_S S'$ and $Y = X' \times_{S'} T$ and
$g$ quasi-compact and quasi-separated, and every abelian sheaf
$\mathcal{F}$ on $T_\etale$ annihilated by $n$
\begin{enumerate}
\item the base change map
$(f')^{-1}R^qg_*\mathcal{F}\to R^qh_*e^{-1}\mathcal{F}$
is injective,
\item if $\mathcal{F} \subset \mathcal{G}$ where $\mathcal{G}$
on $T_\etale$ is annihilated by $n$, then
$$
\Coker\left(
(f')^{-1}R^qg_*\mathcal{F}\to R^qh_*e^{-1}\mathcal{F}
\right)
\subset
\Coker\left(
(f')^{-1}R^qg_*\mathcal{G}\to R^qh_*e^{-1}\mathcal{G}
\right)
$$
\item if in (2) the sheaf $\mathcal{G}$ is an injective sheaf
of $\mathbf{Z}/n\mathbf{Z}$-modules, then
$$
\Coker\left((f')^{-1}R^qg_*\mathcal{F}\to R^qh_*e^{-1}\mathcal{F} \right)
\subset R^qh_*e^{-1}\mathcal{G}
$$
\end{enumerate}
\end{lemma}

\begin{proof}
Choose a short exact sequence
$0 \to \mathcal{F} \to \mathcal{I} \to \mathcal{Q} \to 0$
where $\mathcal{I}$ is an injective sheaf of $\mathbf{Z}/n\mathbf{Z}$-modules.
Consider the induced diagram
$$
\xymatrix{
(f')^{-1}R^{q - 1}g_*\mathcal{I} \ar[d]_{\cong} \ar[r] &
(f')^{-1}R^{q - 1}g_*\mathcal{Q} \ar[d]_{\cong} \ar[r] &
(f')^{-1}R^qg_*\mathcal{F} \ar[d] \ar[r] &
0 \ar[d] \\
R^{q - 1}h_*e^{-1}\mathcal{I} \ar[r] &
R^{q - 1}h_*e^{-1}\mathcal{Q} \ar[r] &
R^qh_*e^{-1}\mathcal{F} \ar[r] &
R^qh_*e^{-1}\mathcal{I}
}
$$
with exact rows. We have the zero in the right upper corner
as $\mathcal{I}$ is injective. The left two vertical arrows are
isomorphisms by $BC(f, n, q - 1)$. We conclude that part (1) holds.
The above also shows that
$$
\Coker\left(
(f')^{-1}R^qg_*\mathcal{F}\to R^qh_*e^{-1}\mathcal{F}
\right)
\subset
R^qh_*e^{-1}\mathcal{I}
$$
hence part (3) holds. To prove (2) choose
$\mathcal{F} \subset \mathcal{G} \subset \mathcal{I}$.
\end{proof}
